\begin{longtable}{@{}l p{12.5cm}@{}}
  \caption{Skills per-instance answer comparison ($-50\% < \Delta \leq -25\%$) for the 9 most divergent items: same question, each model's answer and reasoning.}\label{tab:cmp_skills_neg25} \\
  \toprule
  \endfirsthead
  \toprule
  \endhead
  \multicolumn{2}{@{}l}{\textbf{68221} \quad $\Delta = -50\%$} \\
  \textbf{Question} & You are coaching a graduate student in Data Science who is preparing for a capstone project that requires integrating multiple data sources, implementing machine learning models, and collaborating with domain experts. Using the rehearsal , elaboration , and organization learning strategies, analyze the student's current plan and identify specific strengths, gaps, and recommended actions. In your answer, break down how each strategy can be applied to (1) gather and structure the necessary information, (2) prepare to request and use help from collaborators, and (3) manage resources (time, datasets, compute). Provide concrete examples of rehearsal activities, elaboration techniques, and organizational structures the student should adopt, and explain how these choices address the identified gaps. \newline Give the final answer directly, without reasoning, and do not use the reasoning task such as think\textgreater{}. \\
  \textbf{Reference} & Model answer: \newline Analysis of the current plan (summary and initial appraisal): The student's plan outlines project goals (predictive model for churn), an initial list of data sources (CRM logs, transaction history, product usage telemetry), and a rough timeline for model development and evaluation. Strengths include a clear end goal, awareness of multiple data sources, and intention to consult a domain expert. Gaps are lack of a concrete data integration workflow, no rehearsal activities to validate modeling assumptions, limited strategies for eliciting timely help from collaborators, and an unclear schedule for compute and dataset access. \newline Applying the three strategies to address these gaps: \newline 1) Rehearsal : purpose and concrete activities \newline - Purpose: simulate key tasks to reveal hidden assumptions, surface integration problems early, and build fluency with tools and pipelines before full-scale execution. \newline - Specific rehearsal activities: \newline \textit{[\ldots\ truncated \ldots]} \newline - If cloud GPU quotas are limited, the organization plan will pre-assign heavy experiments to later milestones and reserve smaller instances for iterative work, avoiding idle waiting. \newline Summary (concise): Use rehearsal to surface technical and communication issues early, elaboration to transform raw fields into testable hypotheses and targeted questions, and organization to make workflows reproducible and to manage time, data, and compute. Implement the concrete actions above (synthetic ETL dry-runs, concept maps, feature justifications, a shared data catalog, and a milestone-driven resource plan) to strengthen the student's plan, reduce risk, and make collaborator requests efficient and actionable. \\
  \midrule
  \textbf{Kimi-K2.6} & score 10/10 \\
  & [9-10] The answer demonstrates clear coherence with all strategies, providing highly specific, actionable, and diagnostic examples (e.g., integration fire drills, Conversation Playbook, RACI matrix) tailored to the learner's situation. It explicitly connects actions to identified gaps and resources, showcasing excellent learner-centeredness and immediate feasibility. \\
  \textbf{Innospark-1T} & score 5/10 \\
  & Let me think through this carefully. \newline The question/task: You are coaching a graduate student in Data Science who is preparing for a capstone project that requires integrating multiple data sources, implementing machine learning models, and collaborating with domain experts. Using the rehearsal , elaboration , and organization learning strategies, analyze the student's current plan and ide \newline My reasoning: Coaching a Data Science Capstone Student Using Learning Strategies \newline Analysis of the Student's Current Plan \newline Before applying the three learning strategies, I need to assess what a typical graduate student's plan might look like for this type of capstone project: \newline Typical Student Plan Assessment: \newline \textit{[\ldots\ truncated \ldots]} \newline - Technique: Create decision frameworks with justification \newline - Example: Develop a "compute allocation rationale" document explaining why certain models need GPU vs. \\
  \specialrule{1.2pt}{8pt}{8pt}
  \multicolumn{2}{@{}l}{\textbf{40944} \quad $\Delta = -40\%$} \\
  \textbf{Question} & You are a teaching assistant supporting an instructor during a medical school lecture. During a live clinical reasoning exercise, a student presents a diagnostic plan for a patient with suspected bacterial meningitis that omits immediate collection of blood cultures and delays lumbar puncture until after starting empiric antibiotics. As a teaching assistant responsible for real-time error detection and correction, describe (a) how you would immediately communicate this identified error to the instructor in a concise, actionable way and (b) provide an evidence-based explanation of why this is an error and a recommended correction the instructor can relay to the student. In your explanation, include the clinical rationale, potential patient safety implications, and the suggested sequence of steps the student should have proposed. \newline Give the final answer directly, without reasoning, and do not use the reasoning task such as think\textgreater{}. \\
  \textbf{Reference} & Part (a) --- Immediate concise communication to the instructor: \newline "Instructor, real-time safety concern: the student's plan delays blood cultures and performs lumbar puncture after starting empiric antibiotics. Current guidelines recommend obtaining blood cultures before antibiotics when feasible and performing lumbar puncture promptly unless contraindicated; giving antibiotics first can reduce CSF culture yield and delay targeted therapy. Recommend we pause and clarify sequencing to the student now." \newline Part (b) --- Evidence-based explanation, patient-safety implications, and corrected sequence: \newline The proposed plan is erroneous because starting empiric antibiotics before obtaining blood cultures and before timely lumbar puncture can reduce the diagnostic yield of cerebrospinal fluid (CSF) and blood cultures, hindering identification of the causative organism and its antibiotic susceptibilities. Bacterial meningitis is a time-sensitive, high-risk condition where early appropriate therapy improves outcomes, but accurate microbiologic diagnosis is crucial for narrowing therapy and preventing antimicrobial overuse. \newline Clinical rationale: \newline - Blood cultures should be obtained prior to antibiotics whenever this does not cause harmful delay because they are positive in a substantial proportion of bacterial meningitis cases and guide pathogen identification and susceptibility testing. \newline \textit{[\ldots\ truncated \ldots]} \newline Concise corrective phrasing the instructor can relay to the student: \newline "Good start on the empiric coverage, but remember: obtain blood cultures before antibiotics and perform LP for CSF analysis before antibiotics when it's safe to do so. If LP must be delayed or the patient is unstable, draw blood cultures first and then give antibiotics immediately. This preserves diagnostic yield while ensuring timely therapy." \\
  \midrule
  \textbf{Kimi-K2.6} & score 9/10 \\
  & [9-10] The error analysis is highly accurate and specific, clearly identifying the clinical and diagnostic causes of the error. The correction suggestions are feasible, prioritized, and tailored to the instructor's situation, including a precise corrected sequence and clinical rationale. \\
  \textbf{Innospark-1T} & score 5/10 \\
  & Let me think through this carefully. \newline The question/task: You are a teaching assistant supporting an instructor during a medical school lecture. During a live clinical reasoning exercise, a student presents a diagnostic plan for a patient with suspected bacterial meningitis that omits immediate collection of blood cultures and delays lumbar puncture until after starting empiric antibiotics. As a teaching assistant responsi \newline My reasoning: Teaching Assistant Response: Clinical Reasoning Error Correction \newline Part (a): Immediate Communication to Instructor \newline Slip the instructor a note or use a discreet hand signal, then whisper: \newline \textgreater{} "Critical error in sepsis workup: student omitted blood cultures and proposed LP after antibiotics. Per IDSA guidelines, blood cultures must precede antibiotics; LP timing depends on stability. Recommend pausing to correct this now---significant patient safety issue." \newline \textit{[\ldots\ truncated \ldots]} \newline Suggested Instructor Feedback to Student \newline \textgreater{} "You've identified the right concern, but the sequence needs adjustment. In suspected bacterial meningitis, blood cultures must be obtained before antibiotics---they're our best chance at identifying the organism. The LP is important, but we never delay antibiotics for it. The correct approach is: stabilize, blood cultures, then antibiotics immediately---LP can follow, even if antibiotics have already started. Would you like to walk through this sequence again?" \\
  \specialrule{1.2pt}{8pt}{8pt}
  \multicolumn{2}{@{}l}{\textbf{4250} \quad $\Delta = -40\%$} \\
  \textbf{Question} & As an instructional designer conducting a task analysis for a public health higher-education course, classify the instructional goal 'Students will value and demonstrate respectful communication with diverse patient populations' into Gagn's learning taxonomy (identify which of Gagn's five categories it belongs to). Then, describe the main steps you would include in a goal analysis to determine the components and required abilities for this affective goal at the Receiving level (awareness and willingness to hear). Be specific about what you would identify at each step. \\
  \textbf{Reference} & This instructional goal belongs to Gagn's affective domain category---specifically it aligns with the Affective Learning outcome (attitudes, values, and feelings). Within Gagn's framework, it is an affective-type goal focused on students' willingness to attend to and accept the importance of respectful communication with diverse patients. \newline Main steps for goal analysis targeting Receiving (awareness and willingness to hear): \newline 1) Clarify the terminal goal and its indicators: Define the observable outcomes that show students are at the Receiving level. For this goal I would specify indicators such as: students listen attentively to patients from different backgrounds, acknowledge the importance of cultural respect when prompted, and show openness when introduced to scenarios about diversity. These indicators emphasize awareness and willingness rather than enacted skills. \newline 2) Identify prerequisite attitudes and knowledge: Determine what learners must already be aware of or willing to consider. This includes basic awareness of cultural diversity concepts (e.g., culture, bias, health disparities), prior exposure to the idea that communication affects patient outcomes, and a willingness to attend to new perspectives. At the Receiving level, the prerequisites are primarily receptive---ability to recognize when diverse perspectives are presented and to maintain attention. \newline 3) Break the goal into component sub-skills or sub-attitudes: Decompose the affective goal into smaller, teachable parts appropriate for Receiving. Components include: (a) recognizing respectful vs. disrespectful language, (b) attending to patients' expressions of identity and concern, (c) acknowledging the legitimacy of different cultural health beliefs, and (d) showing openness when presented with unfamiliar cultural practices. Each component is framed as awareness or willingness to hear rather than demonstrated mastery. \newline 4) Determine appropriate learning tasks and stimuli: Select tasks that promote attention and awareness. For Receiving-level outcomes, tasks might include: listening to recorded patient narratives from diverse backgrounds, reading short testimonials about communication experiences in healthcare, participating in guided reflection prompts that ask learners to note their immediate reactions, and attending a live talk from a community representative. These tasks should elicit attention and prompt acceptance of the topic's importance. \newline \textit{[\ldots\ truncated \ldots]} \newline 7) Note potential barriers and motivational strategies: Recognize factors that could impede Receiving (e.g., preconceived resistance, lack of perceived relevance) and include strategies such as establishing relevance to professional practice, using respected role models or patient voices, and creating clear expectations that attention and openness are required. \newline In sum, the goal is classified as affective (Gagn's category for attitudes/values). The goal analysis for the Receiving level focuses on specifying observable indicators of attention and openness, identifying prerequisite awareness, decomposing the goal into receptive components, selecting exposure-based learning tasks, planning simple assessments that demonstrate listening/acknowledgment, sequencing supports to foster receptivity, and addressing barriers to ensure learners are willing to hear and attend to the topic. \\
  \midrule
  \textbf{Kimi-K2.6} & score 9/10 \\
  & [9-10] Correctly identifies Gagn's Attitude domain. Steps are logical, well-sequenced, and use observable verbs (e.g., facing, pausing, noticing). The tree structure is complete, clearly connecting the objective to specific receiving behaviors and prerequisites, appropriately tailored for higher-education learners. \\
  \textbf{Innospark-1T} & score 5/10 \\
  & Let me work through this problem step by step. \newline First, I need to understand what's being asked: As an instructional designer conducting a task analysis for a public health higher-education course, classify the instructional goal 'Students will value and demonstrate respectful communication with diverse patient populations' into Gagn's learning taxonomy (identify which of Gagn's five categories it belongs to). Then, describe the main steps you would include in a goal a \newline Key values given: ,, ,, ,, ,, , \newline Working through the calculation: Analysis and Solution \newline Classification in Gagn's Learning Taxonomy \newline Gagn's Five Categories of Learning Outcomes: \newline \textit{[\ldots\ truncated \ldots]} \newline The instructional goal "Students will value and demonstrate respectful communication with diverse patient populations" is classified as an Attitude outcome in Gagn's learning taxonomy. \newline This gives us the final answer: The instructional goal "Students will value and demonstrate respectful communication with diverse pati \\
  \specialrule{1.2pt}{8pt}{8pt}
  \multicolumn{2}{@{}l}{\textbf{68871} \quad $\Delta = -40\%$} \\
  \textbf{Question} & As a coach helping a higher-education linguistics student reflect on their use of rehearsal , elaboration , and organization strategies, ask the student to describe specific actions they have taken in the past month that demonstrate each of these three strategies, explain how each action helped them manage learning resources or request help from peers/instructors, and identify one concrete change they will make to improve each strategy going forward. \\
  \textbf{Reference} & Over the past month I used distinct actions that illustrate rehearsal, elaboration, and organization strategies. For rehearsal, I scheduled three short, focused review sessions each week where I read aloud key phonetics terms and repeated example transcriptions from lecture notes. This repetition helped commit terminology (e.g., alveolar, palatal, voiced) and transcription conventions to working memory so I could recall them during discussions and quizzes. To manage resources, I set calendar reminders and shared my review schedule with a study partner so we could coordinate brief check-ins; when I struggled with a concept I messaged the partner to request a quick clarification. For elaboration, I created concept maps linking morphology concepts to real language examples and wrote short explanatory paragraphs connecting theoretical models to data from a course corpus. These elaborative activities forced me to rephrase definitions in my own words and integrate new ideas with prior knowledge, which made it easier to formulate targeted questions to ask the instructor during office hours. I also annotated articles with margin notes that connected empirical findings to seminar topics, and when I needed deeper feedback I emailed the instructor specific excerpts with my interpretations. For organization, I grouped my notes by topic (phonology, morphology, syntax) and maintained a shared folder with labeled subfolders for readings, datasets, and assignments. I used consistent file names and a single reference list document to reduce time spent searching for sources. When collaborative work was required, I created a shared checklist for task division so team members knew what to do and when to ask for help. Concrete changes I will make: (1) Rehearsal --- I will incorporate spaced retrieval by increasing intervals between review sessions and using self-quizzing flashcards rather than passive rereading to improve long-term retention. (2) Elaboration --- I will add brief peer-teaching sessions where I explain a concept to a classmate weekly, which will reveal gaps in my understanding and prompt more specific help requests when needed. (3) Organization --- I will adopt a digital note-taking template that includes a 'questions' section for each topic so that when I encounter unclear points I can compile them and request focused feedback from peers or the instructor during scheduled meeting times. \\
  \midrule
  \textbf{Kimi-K2.6} & score 9/10 \\
  & [9-10] The coaching questions are highly specific, learner-centered, and tailored to a linguistics student with relevant examples. They clearly address all four strategies (rehearsal/demonstration, elaboration, organization, resource management) and prompt the learner to derive a concrete action plan, fitting the 9-10 band. \\
  \textbf{Innospark-1T} & score 5/10 \\
  & Let me think through this carefully. \newline The question/task: As a coach helping a higher-education linguistics student reflect on their use of rehearsal , elaboration , and organization strategies, ask the student to describe specific actions they have taken in the past month that demonstrate each of these three strategies, explain how each action helped them ma \newline My reasoning: Coaching Reflection on Learning Strategies \newline Thinking Process \newline This request asks me to act as a coach for a linguistics student, guiding them through reflection on three specific learning strategies: \newline 1. Rehearsal - Repeating or practicing information to aid retention \newline \textit{[\ldots\ truncated \ldots]} \newline - Organization creates structure that makes information accessible \newline Your action plan should include specific, measurable steps for each strategy. Would you like to share your responses to these questions, or do you need clarification on any of them?" \\
  \specialrule{1.2pt}{8pt}{8pt}
  \multicolumn{2}{@{}l}{\textbf{68996} \quad $\Delta = -40\%$} \\
  \textbf{Question} & You are coaching a university-level ensemble preparing for a public performance. Using the rehearsal strategies of retrieval practice , elaboration , and organization , analyze the ensemble's current preparation plan and identify specific modifications to improve information gathering and help-seeking behaviors. In your response, break down which deficiencies each strategy would address, propose concrete rehearsal activities or checkpoints that apply each strategy, and explain how you would prompt the ensemble members to seek appropriate external resources or peer assistance at each stage. \\
  \textbf{Reference} & Current preparation plan analysis and proposed modifications: \newline 1) Retrieval practice : \newline - Identified deficiencies: The ensemble currently relies on passive run-throughs during rehearsals, which encourages familiarity but not robust memory or quick error detection. Members often read parts rather than retrieve from memory, leading to insecure entrances and inconsistent cueing. \newline - Concrete modifications: Introduce short, frequent closed-section retrieval drills where players/singers perform small excerpts from memory or from reduced visual cues (e.g., score covered except measures to be played). Implement random cold-starts where a conductor cues an instrument to begin without prior countdown to simulate performance conditions. Schedule brief end-of-session recall checks where each section must play a key motif or cue sequence from memory. \newline - Help-seeking prompts: After a retrieval drill, ask each section leader to report which measures or transitions caused hesitation and to request targeted help (e.g., rhythmic counting strategies from the percussion coach, harmonic clarification from the music theory tutor). Encourage members to use recorded rehearsals to compare their retrieval attempts with earlier takes and to share timestamps where problems recur. \newline 2) Elaboration : \newline \textit{[\ldots\ truncated \ldots]} \newline - What top two items go on the rehearsal map today, and what external assistance do you need to clear them before the performance rehearsal? \newline Overall, by diagnosing passive rehearsal habits and replacing them with coordinated retrieval, elaboration, and organization activities, the ensemble will gather higher-quality diagnostic information and adopt clearer, faster help-seeking routines that directly target performance weaknesses. \\
  \midrule
  \textbf{Kimi-K2.6} & score 9/10 \\
  & [9-10] The response clearly addresses all strategies (retrieval/demonstration, elaboration, organization, resource management) with highly specific, diagnostic prompts and concrete action plans. It effectively connects help-seeking to resources (librarian, principals, peer partners) and overcomes obstacles via the P1-P3 matrix and structured workflows, demonstrating excellent learner-centeredness and feasibility. \\
  \textbf{Innospark-1T} & score 5/10 \\
  & Let me think through this carefully. \newline The question/task: You are coaching a university-level ensemble preparing for a public performance. Using the rehearsal strategies of retrieval practice , elaboration , and organization , analyze the ensemble's current preparation plan and identify specific modifications to improve information gathering and help-seeking behaviors. In your response, break down whi \newline My reasoning: \textless{}analysis\textgreater{} \newline Let me analyze this educational coaching scenario systematically: \newline 1. Understanding the Core Framework: \newline - Retrieval practice : Active recall of information from memory without cues \newline \textit{[\ldots\ truncated \ldots]} \newline - Quality of questions asked during rehearsals \newline By systematically applying retrieval practice, elaboration, and organization strategies, the ensemble will develop deeper musical understanding, stronger performance capabilities, and more effective collaborative learning behaviors essential for successful public performance. \\
  \specialrule{1.2pt}{8pt}{8pt}
  \multicolumn{2}{@{}l}{\textbf{105002} \quad $\Delta = -40\%$} \\
  \textbf{Question} & You are an evaluator analyzing a student's incorrect answer on an exam question about implementing a thread-safe cache in a concurrent web server. The student chose to synchronize all cache accesses using a single global lock, and their incorrectness report stated: "Using one global lock ensures consistency, so it's correct." Based on the assessment data (correct/incorrect, error type, and inferred understanding level), analyze the likely misconception(s) behind this answer. In your response, (1) identify and break down the underlying conceptual errors and their probable causes, (2) compare the global-lock approach with at least two alternative concurrency-control strategies (e.g., fine-grained locks, lock-free/wait-free structures, read-write locks, or concurrent hash maps), explaining trade-offs in correctness, performance, and complexity, (3) give a personalized explanation aimed at correcting the student's misconceptions (include conceptual clarifications and a short illustrative example), and (4) recommend a targeted re-learning strategy with concrete activities or exercises (including one small coding task and one reading or visualization exercise) to help the student internalize correct concurrency design principles. Your analysis should demonstrate how each part of your feedback addresses the specific misconception that "one global lock is sufficient." \\
  \textbf{Reference} & Likely misconceptions and their causes \newline 1) Overgeneralization of correctness from simplicity: The student believes that because a single global lock prevents concurrent conflicting accesses, it therefore is an acceptable or optimal solution. This conflates correctness (mutual exclusion ensures safety) with adequacy for system requirements (performance, scalability, and responsiveness). Cause: limited exposure to trade-offs in concurrent design and insufficient experience measuring performance under contention. \newline 2) Failure to distinguish safety and liveness/performance: The student treats "consistency" as the only concern, overlooking latency, throughput, and blocking behavior (e.g., priority inversion, convoying). Cause: conceptual gap between functional correctness and non-functional properties of concurrent systems. \newline 3) Lack of understanding of workload characteristics and contention: The student assumes uniform access patterns; they don't recognize that read-heavy or partitionable workloads can be optimized without sacrificing correctness. Cause: absence of workload analysis practice and profiling experience. \newline 4) Underappreciation of composability and modularity costs: Using a single lock can make composing subsystems harder (coarse-grained locking increases coupling and risk of deadlock when combined with other locks). Cause: little practice with integrating concurrent components. \newline Comparison of global lock vs. alternatives (trade-offs) \newline \textit{[\ldots\ truncated \ldots]} \newline Summary (concise takeaways for the student) \newline - A single global lock gives correctness but often fails to meet performance and scalability requirements. Concurrency design requires analyzing workload and choosing an approach (read-write locks, sharding, or concurrent data structures) that balances safety, performance, and complexity. Use experiments and profiling to validate design choices rather than assuming a simple solution is sufficient. \\
  \midrule
  \textbf{Kimi-K2.6} & score 10/10 \\
  & [9-10] The response accurately identifies and classifies the misconceptions (safety-liveness conflation, lock-granularity blindness, domain decontextualization) with clear probable causes. It compares the global lock with two alternatives (Read-Write Lock, Fine-Grained Locks) across correctness, performance, and complexity. The personalized explanation uses a relatable analogy and concrete pseudo-code to directly correct the misconception. The re-learning strategy is highly actionable, featuring a specific coding benchmark, a targeted reading/visualization exercise, and a reflection prompt that serves as a method to confirm understanding. \\
  \textbf{Innospark-1T} & score 6/10 \\
  & Let me think through this carefully. \newline The question/task: You are an evaluator analyzing a student's incorrect answer on an exam question about implementing a thread-safe cache in a concurrent web server. The student chose to synchronize all cache accesses using a single global lock, and their incorrectness report stated: "Using one global lock ensures consistency, so it's correct." Based on the assessme \newline My reasoning: Analysis of Student Misconception: Thread-Safe Cache Implementation \newline Part 1: Identifying and Breaking Down Conceptual Errors \newline Primary Misconception: "Correctness = Consistency" \newline The student conflates two distinct concepts: \newline \textit{[\ldots\ truncated \ldots]} \newline This comprehensive approach addresses the core misconception by: \newline 1. Providing hands-on experience with performance implications \\
  \specialrule{1.2pt}{8pt}{8pt}
  \multicolumn{2}{@{}l}{\textbf{100140} \quad $\Delta = -30\%$} \\
  \textbf{Question} & You are evaluating a student's environmental science project report about local freshwater quality. The rubric requires that you determine the student's achievement level by scoring each rubric criterion, assigning partial scores where appropriate, and identifying common error types that affected the score. Describe, in detail, the step-by-step process you would use to: 1) score each rubric criterion and calculate partial/total scores; 2) classify and record the types of errors or misconceptions found in the student's data collection, analysis, and conclusions; and 3) use those scores and error classifications to justify an overall achievement level for the summative assessment. Be specific about what evidence you would cite from the student's report for each step and how that evidence maps to rubric descriptors. \\
  \textbf{Reference} & Step 1 --- Score each rubric criterion and calculate partial/total scores \newline 1. Gather materials and rubric: I begin by having the student's project report, any attached raw data (e.g., water quality measurements), field notes, and the rubric with clearly defined performance descriptors for each criterion (for example: Methodology, Data Quality, Analysis \& Interpretation, Use of Evidence, Communication). \newline 2. Read the report holistically: I read the entire report once to understand its scope and main claims. This helps me see how parts connect before assigning scores. \newline 3. Score criterion-by-criterion using descriptors: For each rubric criterion, I compare the student's evidence to the descriptor language. If a criterion uses a 0--4 scale, I identify the descriptor that best matches the student's performance. For example: \newline - Methodology: If the student provided a clear sampling plan with specified locations, consistent sampling times, and described instruments and calibration steps, I select the descriptor for "proficient" (e.g., 3 out of 4). If the plan is incomplete (e.g., locations listed but no replication or missing calibration details), I assign a partial score (e.g., 2 out of 4). \newline - Data Quality: If measurements include units, repeated trials, and documented anomalies, that supports a higher score; if values are missing, units inconsistent, or obvious transcription errors appear, I reduce the score accordingly. \newline \textit{[\ldots\ truncated \ldots]} \newline Final summary statement I would record as the evaluator: Overall achievement: Proficient. The student demonstrates a clear understanding of freshwater monitoring concepts and presents appropriate basic analyses, but methodological shortcomings (single samples, missing calibration), a transcription/calculation error that affects reported summary statistics, and a tendency to overinterpret correlational data limit the reliability of conclusions. Specific evidence for these determinations is cited above by page and table/figure. \newline This process ensures each numeric score and the chosen achievement level are traceable to explicit evidence and that error classifications explain why certain rubric criteria received partial credit. \\
  \midrule
  \textbf{Kimi-K2.6} & score 10/10 \\
  & [9-10] The student's response thoroughly and accurately addresses all aspects of the prompt. The classification of error types is highly detailed and well-categorized (Procedural, Analytical, Conceptual, Inferential, Communication), perfectly aligning with the rubric's requirement for classifying incorrect answer types. The validity of partial scoring is explicitly justified through component decomposition with clear examples. The basis for the final achievement level is logically connected to the item results, utilizing a structured summative justification statement that directly maps evidence and error patterns to rubric descriptors. \\
  \textbf{Innospark-1T} & score 7/10 \\
  & Let me think through this carefully. \newline The question/task: You are evaluating a student's environmental science project report about local freshwater quality. The rubric requires that you determine the student's achievement level by scoring each rubric criterion, assigning partial scores where appropriate, and identifying common error types that affected the score. Describe, in detail, the step-by-step process you would use to: 1) score each ru \newline My reasoning: Step-by-Step Evaluation Process for Freshwater Quality Project Report \newline Step 1: Scoring Each Rubric Criterion and Calculating Partial/Total Scores \newline Preparation Phase \newline - Review the complete rubric beforehand to internalize performance level descriptors \newline \textit{[\ldots\ truncated \ldots]} \newline - Achievement level justified with specific evidence \newline - Feedback prepared that addresses both strengths and areas for improvement \\
  \specialrule{1.2pt}{8pt}{8pt}
  \multicolumn{2}{@{}l}{\textbf{23467} \quad $\Delta = -30\%$} \\
  \textbf{Question} & You are designing a transition-focused lesson for high-school students with intellectual disabilities aimed at improving independent community transportation planning (e.g., planning a bus trip from school to a job interview location). Using Gagn's Nine Events of Instruction and the ABCD model for writing learning objectives, create a complete lesson plan outline that aligns with a single, specific ABCD-formatted learning objective. Your response should: (1) state the learning objective in ABCD format including audience, behavior, condition, and degree/criteria and clearly specify the instructional content; (2) justify your choice of instructional media and materials tied to the objective and learners' needs; (3) map each of Gagn's nine events to concrete classroom activities or materials (including scaffolds, prompts, formative checks, and ways to fade supports) that will enable learners to achieve the objective; (4) describe assessment methods and the exact criteria you will use to determine mastery (how criterion maps to the objective's degree); and (5) explain how you will adapt the lesson for two levels of learner variability (one lower-support need and one higher-support need) while preserving alignment with the objective. Ensure the plan demonstrates synthesis, design judgment, and clear alignment between objective, media, instructional events, and assessment. \newline Give the final answer directly, without reasoning, and do not use the reasoning task such as think\textgreater{}. \\
  \textbf{Reference} & Learning objective (ABCD): Given a simplified transit map, bus schedule, and a written scenario (Condition), five high-school students with mild-to-moderate intellectual disabilities (Audience) will independently plan and record a bus trip that departs school at least 30 minutes before a scheduled 9:00 AM job interview, includes two correct transfers if required, and lists the expected arrival time within 10 minutes (Degree) by producing a written trip plan and verbally explaining each step (Behavior). Instructional content: interpreting a simplified transit map and schedule, identifying departure times, calculating trip duration (including transfer times and walking time), sequencing steps for boarding and transfers, and producing a written/verbal trip plan for a real-world scenario. \newline Justification of media and materials: I select a multimodal set: printed simplified transit maps with large icons and color-coded routes, laminated bus schedules with emphasized key times, step-by-step trip planning checklist cards, a visual timetable template for learners to fill, tablet-based interactive route simulator (visual animation of stops), and real-world props (mock bus pass, play phone, printed interview address). These media support visual processing, reduce cognitive load by simplifying extraneous detail, and allow repeated practice (simulation) and multimodal expression (written + verbal). The interactive simulator provides safe, repeatable trials and immediate feedback; printed templates scaffold production of the written plan. Materials are chosen to match learners' strengths in visual aided learning and to permit fading of supports. \newline Mapping Gagn's Nine Events to concrete activities and scaffolds: \newline 1) Gain attention: Begin with a short, two-minute animated simulation on the tablet that shows a character successfully reaching an interview on time by taking the bus. Prompt: "Watch how Maya gets there on time---what steps does she take?" (Purpose: activate interest and set a real-world goal.) \newline 2) Inform learners of objectives: Display the ABCD objective in student-friendly language and read it aloud: "Today you will plan and write a bus trip so you arrive 30 minutes early to a 9:00 interview." Provide a one-page checklist of success criteria (includes degree: arrival within 10 minutes, independently listing steps). \newline 3) Stimulate recall of prior learning: Use a quick two-question warm-up asking students to name parts of a bus schedule and to identify a bus stop icon on the map. Provide prompts for those who need it (picture cue cards). This activates existing knowledge of maps and schedules. \newline \textit{[\ldots\ truncated \ldots]} \newline - Higher-support need learner (more assistance): Increase scaffolds and break tasks into smaller subskills while keeping the same overall objective. Steps: (1) Pre-teach time addition using manipulatives (minute cards) and practice matching departure times to arrival times; (2) provide a partially-completed template where departure time and final arrival are highlighted and students only need to fill middle steps initially; (3) permit one additional prompting cue during the summative verbal explanation (so up to two prompts total); (4) allow the learner to use the simulator in guided mode where the teacher pauses at each transfer decision. To ensure alignment, the final summative expectation remains the same ABCD objective, but degree is implemented by allowing an accommodated path to demonstrate the same behaviors (e.g., the student still must identify correct route and arrival within 10 minutes, but the teacher documents prompts used). Over time fade supports: progressively remove worked examples, reduce hint frequency, and shift from teacher-led simulator to independent runs. \newline Synthesis and design judgment summary: This lesson plan uses the ABCD objective to anchor every instructional decision: chosen media (visual maps, simulator, templates) reduce cognitive load and provide multimodal practice; Gagn's nine events structure learning from attention to transfer; scaffolds and graduated prompts target skill components; and assessments are directly mapped to objective criteria (arrival time tolerance, correct transfers, independent verbalization). Adaptations are designed to maintain fidelity to the objective while providing differentiated pathways to mastery, and fading strategies ensure movement toward independent performance. \\
  \midrule
  \textbf{Kimi-K2.6} & score 10/10 \\
  & [9-10] The response demonstrates comprehensive alignment between the ABCD objective, instructional media, Gagn's events, and assessment. Differentiation is organically integrated with clear fading strategies, and evidence-based practices (e.g., task analysis, least-to-most prompting, concrete-representational scaffolds) are used to ensure accessibility and active participation for diverse learners. \\
  \textbf{Innospark-1T} & score 7/10 \\
  & Let me think through this carefully. \newline The question/task: You are designing a transition-focused lesson for high-school students with intellectual disabilities aimed at improving independent community transportation planning (e.g., planning a bus trip from school to a job interview location). Using Gagn's Nine Events of Instruction and the ABCD model for writing learning objectives, create a complete lesson plan outline that alig \newline My reasoning: Lesson Plan: Independent Community Transportation Planning \newline 1. Learning Objective (ABCD Format) \newline Given a scenario card describing a job interview location (condition), the student (audience) will independently plan a bus trip from school to the destination (behavior) with 100\% accuracy in identifying the correct bus route, stop, transfer points (if needed), arrival time, and fare requirements (degree/criteria). Instructional Content \newline Core Content: \newline \textit{[\ldots\ truncated \ldots]} \newline | Adaptations | Two levels of support | Preserves objective alignment | \newline This lesson plan demonstrates that effective instruction for students with intellectual disabilities requires careful alignment between objectives, materials, instructional strategies, and assessment---while remaining flexible enough to meet individual learner needs within a unified instructional framework. \\
  \specialrule{1.2pt}{8pt}{8pt}
  \multicolumn{2}{@{}l}{\textbf{3657} \quad $\Delta = -30\%$} \\
  \textbf{Question} & You are an instructional designer performing a task analysis for a higher-education mathematics course. The instructional goal is: "Students will be able to formulate and evaluate rigorous proofs of multivariable calculus theorems involving gradient, divergence, and curl, and determine appropriate theorems to apply in novel contexts." Using Gagn's learning taxonomy, classify this instructional goal into the appropriate learning domain and specific category (e.g., verbal information, intellectual skills---divisible into rules, discriminations, concrete concepts, higher-order rules, etc.). Then, provide a detailed main-step analysis (task analysis) that breaks this goal into constituent performance components and subskills, sequences them from simple to complex, and justifies each sequencing and inclusion based on cognitive demands. Finally, evaluate at least two alternative task decompositions (e.g., emphasizing procedural computation first vs. emphasizing conceptual proofs first), justify which decomposition best aligns with the goal and the target learners (advanced undergraduate/early graduate mathematics students), and recommend instructional entry conditions and prerequisite objectives. Provide clear, evidence-based reasoning for your design choices using concepts from Gagn's taxonomy and higher-order cognitive processes (evaluating, creating). \newline Give the final answer directly, without reasoning, and do not use the reasoning task such as think\textgreater{}. \\
  \textbf{Reference} & Model Answer: \newline Classification according to Gagn's taxonomy \newline - Learning domain: Cognitive. \newline - Specific category: Intellectual skills---higher-order rules (problem-solving and theorem construction) and discriminations. Rationale: The goal requires students to recognize when particular theorems (e.g., Green's, Divergence, Stokes') apply (discrimination), to synthesize conditions and constructs into rigorous proofs (higher-order rules), and to create and evaluate novel proofs in unfamiliar contexts. Verbal information (definitions, theorem statements) and motor skills (notation manipulation) are supporting but not the central category. \newline Main-step (task) analysis: decomposition into constituent components and sequencing \newline Overview: Decompose into foundational knowledge items (verbal information), then conceptual discriminations, then basic procedures, then integrated proof-construction skills, culminating in evaluation and creation of novel proofs. \newline \textit{[\ldots\ truncated \ldots]} \newline Final justification using Gagn's taxonomy and higher-order cognitive processes \newline - Classifying the goal as intellectual skills---higher-order rules emphasizes that students must form, apply, and critique rule systems (theorems) in novel contexts. The task analysis orders learning from core verbal information and discriminations through procedural fluency to synthesis and evaluation, which matches Gagn's hierarchy (simple verbal/conceptual elements discriminations higher-order rule application). The recommended integrated decomposition supports the Evaluating-level demands by prioritizing conceptual clarity and discrimination while ensuring procedural competence is sufficient to execute and check proofs. This design fosters students' ability to judge validity, generalize results, and create new proofs---fulfilling the stated instructional goal. \\
  \midrule
  \textbf{Kimi-K2.6} & score 10/10 \\
  & [9-10] The student correctly identifies the Gagn domain (Intellectual Skills---Higher-Order Rules/Cognitive Strategies) and aligns it perfectly with the advanced learner level. Main steps are clearly described using observable verbs (recall, distinguish, classify, compute, analyze, assemble, critique), logically sequenced from simple to complex, and form a complete, well-structured tree. All criteria are fully met. \\
  \textbf{Innospark-1T} & score 7/10 \\
  & Let me think through this carefully. \newline The question/task: You are an instructional designer performing a task analysis for a higher-education mathematics course. The instructional goal is: "Students will be able to formulate and evaluate rigorous proofs of multivariable calculus theorems involving gradient, divergence, and curl, and determine appropriate theorems to apply in novel contexts." Using Gagn's \newline My reasoning: Task Analysis for Multivariable Calculus Proofs \newline Classification Using Gagn's Learning Taxonomy \newline Domain Analysis \newline This instructional goal spans multiple domains in Gagn's taxonomy: \newline \textit{[\ldots\ truncated \ldots]} \newline - [x] Listed prerequisite objectives with Gagn categories \newline - [x] Provided evidence-based reasoning throughout \\
  \specialrule{1.2pt}{8pt}{8pt}
  \bottomrule
\end{longtable}